\section{TD1}


\begin{td-exo}
    On a
    \begin{equation*}
        \begin{cases}
            c_1 &\colon x_1 + x_2 = x_3 \\
            c_2 &\colon x_1 < x_2
        \end{cases}
    \end{equation*}
    avec
    \begin{equation*}
        \begin{cases}
            D(x_1) &= \left\{1, 2, 3, 4, 5\right\} \\
            D(x_2) &= \left\{1, 3, 5\right\} \\
            D(x_3) &= \left\{1, 3, 4, 7\right\}.
        \end{cases}
    \end{equation*}

    On trouve
    \begin{equation*}
        c_2 \colon (x_1, 5), (x_2, 1) \quad\text{et}\quad c_1 \colon (x_3, 1), (x_3, 3), (x_1, 3).
    \end{equation*}

    Après \(AC\), on a donc trouvé
    \begin{equation*}
        \begin{cases}
            D(x_1) &= \left\{1, 2, 4\right\} \\
            D(x_2) &= \left\{3, 5\right\} \\
            D(x_3) &= \left\{4, 7\right\}.
        \end{cases}
    \end{equation*}
\end{td-exo}

\begin{td-exo}
    On a
    \begin{equation*}
        \begin{cases}
            c_1 &\colon 2x_1 = x_3 \\
            c_2 &\colon x_1 < x_2 \\
            c_3 &\colon x_3 + x_4 = 10 \\
            c_4 &\colon x_3 < x_2
        \end{cases}
    \end{equation*}
    avec
    \begin{equation*}
        \forall i \in \left\{1, 4\right\},\quad D(x_i) = \left\{1, \ldots, 8\right\}.
    \end{equation*}

    On trouve
    \begin{equation*}
        \begin{aligned}
            c_1 &\colon (x_3, \left\{1, 3, 5, 7\right\}), (x_1, \left\{5, 6, 7, 8\right\}) \\
            c_2 &\colon (x_2, 1) \\
            c_3 &\colon (x_4, \left\{1, 3, 5, 7\right\}).
            c_4 &\colon (x_3, 8), (x_2, 2).
        \end{aligned}
    \end{equation*}

    Après \(AC\), on a donc trouvé
    \begin{equation*}
        \begin{cases}
            D(x_1) &= \left\{1, 2, 3\right\} \\
            D(x_2) &= \left\{3, \ldots, 7\right\} \\
            D(x_3) &= \left\{2, 4, 6\right\} \\
            D(x_4) &= \left\{4, 6, 8\right\}.
        \end{cases}
    \end{equation*}
\end{td-exo}

\begin{td-exo}
    On a
    \begin{equation*}
        \begin{cases}
            c_1 &\colon x_1 + x_4 = 10 - 5x_5 \\
            c_2 &\colon |x_1 - x_2| = |x_3 - x_4| \\
            c_3 &\colon x_3 = 2x_2
        \end{cases}
    \end{equation*}
    avec
    \begin{equation*}
        \begin{cases}
            D(x_1) &= \left\{1, 2, 7,  8\right\} \\
            D(x_2) &= \left\{1, 2, 7,  8\right\} \\
            D(x_3) &= \left\{1, 2, 3, 4\right\} \\
            D(x_4) &= \left\{1, 2, 3, 4\right\} \\
            D(x_4) &= \left\{0, 1\right\}.
        \end{cases}
    \end{equation*}

    On trouve
    \begin{equation*}
        \begin{aligned}
            c_1 &\colon (x_1, \left\{5, 6, 7, 8\right\}), (x_3, \left\{1, 3, 5, 7\right\})
            c_2 &\colon (x_2, 1) \\
            c_3 &\colon (x_4, \left\{1, 3, 5, 7\right\}).
            c_4 &\colon (x_3, 8), (x_2, 2).
        \end{aligned}
    \end{equation*}

    Après \(AC\), on a donc trouvé
    \begin{equation*}
        \begin{cases}
            D(x_1) &= \left\{1, 2\right\} \\
            D(x_2) &= \left\{1, 2\right\} \\
            D(x_3) &= \left\{2, 4\right\} \\
            D(x_4) &= \left\{3, 4\right\} \\
            D(x_5) &= \left\{1\right\}.
        \end{cases}
    \end{equation*}
\end{td-exo}

\begin{td-exo}
    Soit le réseau de contraintes \(N = (X, D, C)\) où \(X = \left\{ X_1,\ldots X_n\right\}\) avec pour domaines :
    \begin{equation*}
        D = 
        \begin{cases}
            D(X_1) &= \left\{0, 1, 2, 3, 6\right\} \\
            D(X_2) &= \left\{1, 2, 3, 4\right\} \\
            D(X_3) &= \left\{0, 1, 3, 4, 5\right\} \\
            D(X_4) &= \left\{3, 4, 6, 10, 15\right\} \\
            D(X_5) &= \left\{0, 3, 4\right\}
        \end{cases}
    \end{equation*}
    et pour contraintes :
    \begin{equation*}
        C = 
        \begin{cases}
            X_1 + X_2 + X_3 = X_4 \\
            |X_1 - X_3| = X_5 \\
            |X_1 - X_2| = 2 \\
            X_3 + X_5 \neq X_1 + X_2 + X_4
        \end{cases}
    \end{equation*}
    Pour gagner du temps, on admettra que \(N\) est arc cohérent.
    \begin{enumerate}
        \item Appliquer l'algorithme MAC au réseau de contraintes \(N\).
        On utilisera l'heuristique dom/deg pour le choix des variables de l'ordre lexicographique pour le choix des valeurs.
        Si plusieurs variables ont le m\^eme ratio dom/deg, MAC choisira en priorité la variable \(X_i\) avec l'indice \(i\) le plus petit.
    \end{enumerate}
\end{td-exo}

\begin{td-exo}
    Soit le réseau de contraintes \(N = (X, D, C)\) où \(X = \left\{X_1,\ldots,X_4\right\}\) avec pour domaines :
    \begin{equation*}
        D = 
        \begin{cases}
            D(X_1) &= \left\{1, 2, 7, 8\right\} \\
            D(X_2) &= \left\{1, 2, 7, 8\right\} \\
            D(X_3) &= \left\{1, 2, 3, 4\right\} \\
            D(X_4) &= \left\{1, 2, 3, 4\right\}
        \end{cases}
    \end{equation*}
    et pour contraintes :
    \begin{equation*}
        C = 
        \begin{cases}
            X_1 + X_4 = 5 \\
            |X_1 - X_2| = |X_3 - X_4| \\
            X_3 = 2 X_2
        \end{cases}
    \end{equation*}
    \begin{enumerate}
        \item Appliquer BC sur \(N\).

        \item Appliquer SAC sur \(N\).

        \item Combien comporte de variables le plus petit réseau de contraintes binaires normalisé (c'est-à-dire tel que deux contraintes ne peuvent pas avoir la même portée) qui est SAC, mais n'a pas de solution ? Justifiez votre réponse.

        \item On peut définir la propriété SBC de façon analogue à SAC, en remplaçant AC par BC dans la définition. Comparez les propriétés AC, SAC, BC, SBC.\\
        Un réseau de contraintes binaires est \(k\)-quasi arborescent s'il est normalisé et qu'il existe \(k\) variables dont la suppression (ainsi que la suppression des contraintes incidentes) produit un réseau arborescent.

        \item Un réseau \(1\)-quasi arborescent qui est SAC admet-il toujours une solution ? Justifiez votre réponse.

        \item Soit \(k\geq 1\). Peut-on résoudre en temps polynomial les réseaux de contraintes \(k\)-quasi arborescents ? Justifiez votre réponse.
    \end{enumerate}
\end{td-exo}